\begin{PSbox}[unbreakable]{obj}{Learning Objectives}

After completing this module, students will be able to:

\begin{enumerate}
\def\labelenumi{\arabic{enumi}.}
\tightlist
\item
  Identify which probability distribution models a given engineering scenario
\item
  State the $\frac{\mathrm{PMF}}{\mathrm{PDF}}$, CDF, mean, and variance for all major distributions
\item
  Explain the physical interpretation and engineering context for each distribution
\item
  Determine when to use (and when NOT to use) each distribution
\item
  Implement distributions in MATLAB for simulation
\item
  Recognize distribution patterns from data characteristics
\end{enumerate}

\end{PSbox}

\PSsection{3.1}{Introduction: Choosing the Right Distribution}

\PSsubsection{The Engineering Challenge}

When modeling a random phenomenon, the critical question is:

\begin{PSquote}

``Which probability distribution best describes this uncertain quantity?''

\end{PSquote}

Choosing the wrong distribution leads to incorrect predictions, improper designs, and system failures.

\PSsubsection{How to Choose}

The distribution should match the \textbf{physical mechanism} generating randomness:

\begin{itemize}
\tightlist
\item
  Counting successes \ensuremath{\to} Binomial
\item
  Waiting for events \ensuremath{\to} Exponential/Geometric
\item
  Sum of many small effects \ensuremath{\to} Gaussian
\item
  Signal envelope in fading \ensuremath{\to} Rayleigh
\end{itemize}

This module provides the tools to make this choice correctly.

\PSsection{3.2}{Bernoulli Distribution}

\PSsubsection{Physical Interpretation}

Models a \textbf{single trial} with two possible outcomes: ``success'' (1) or ``failure'' (0).

\PSsubsection{Engineering Context}

Any binary random outcome:

\begin{itemize}
\tightlist
\item
  One bit transmitted: error or correct
\item
  One component tested: pass or fail
\item
  One packet sent: delivered or lost
\item
  One detection attempt: target present or absent
\end{itemize}

\begin{PSbox}[unbreakable]{def}{Mathematical Definition}

$X \sim \mathrm{Bernoulli}(p)$, where $p$ = probability of success.

\end{PSbox}

\PSsubsection{PMF}

\PSdisplay{p_X(x) = \begin{cases} p & x = 1 \\ 1-p = q & x = 0 \end{cases}}

Compact form: $p_{X}(x) = p^{x}(1-p)^{1-x}$ for $x \in \{0,\, 1\}$

\PSsubsection{CDF}

\PSdisplay{F_X(x) = \begin{cases} 0 & x < 0 \\ 1-p & 0 \leq x < 1 \\ 1 & x \geq 1 \end{cases}}

\PSsubsection{Mean}

\PSdisplay{E[X] = p}

\PSsubsection{Variance}

\PSdisplay{\text{Var}(X) = p(1-p)}

Maximum variance at $p = 0.5$ (maximum uncertainty).

\PSsubsection{When to Use}

\PSyes{} Single binary trial with fixed probability\PSbr{}\PSyes{} Building block for more complex distributions\PSbr{}\PSyes{} Modeling $\frac{ON}{\mathrm{OFF}}$ states, pass/fail outcomes

\PSsubsection{When NOT to Use}

\PSno{} Multiple trials (use Binomial)\PSbr{}\PSno{} Outcome is not binary\PSbr{}\PSno{} Probability changes between trials

\PSsubsection{MATLAB Implementation}

\tcbset{PScodemode/.style={unbreakable}}

\begin{Shaded}
\begin{Highlighting}[]
\CommentTok{\% Generate Bernoulli random variables}
\VariableTok{p} \OperatorTok{=} \FloatTok{0.1}\OperatorTok{;}          \CommentTok{\% Bit error probability}
\VariableTok{N} \OperatorTok{=} \FloatTok{10000}\OperatorTok{;}
\VariableTok{X} \OperatorTok{=} \VariableTok{binornd}\NormalTok{(}\FloatTok{1}\OperatorTok{,} \VariableTok{p}\OperatorTok{,} \FloatTok{1}\OperatorTok{,} \VariableTok{N}\NormalTok{)}\OperatorTok{;}  \CommentTok{\% or: X = rand(1,N) \textless{} p;}

\CommentTok{\% Verify}
\VariableTok{fprintf}\NormalTok{(}\SpecialStringTok{\textquotesingle{}Mean: \%.4f (theory: \%.4f)\textbackslash{}n\textquotesingle{}}\OperatorTok{,} \VariableTok{mean}\NormalTok{(}\VariableTok{X}\NormalTok{)}\OperatorTok{,} \VariableTok{p}\NormalTok{)}\OperatorTok{;}
\VariableTok{fprintf}\NormalTok{(}\SpecialStringTok{\textquotesingle{}Var:  \%.4f (theory: \%.4f)\textbackslash{}n\textquotesingle{}}\OperatorTok{,} \VariableTok{var}\NormalTok{(}\VariableTok{X}\NormalTok{)}\OperatorTok{,} \VariableTok{p}\OperatorTok{*}\NormalTok{(}\FloatTok{1}\OperatorTok{{-}}\VariableTok{p}\NormalTok{))}\OperatorTok{;}
\end{Highlighting}
\end{Shaded}

\PSsection{3.3}{Binomial Distribution}

\PSsubsection{Physical Interpretation}

Models the \textbf{number of successes} in $n$ \textbf{independent} trials, each with the same success probability $p$.

\PSsubsection{Engineering Context}

\begin{itemize}
\tightlist
\item
  Number of bit errors in $n$ transmitted bits
\item
  Number of defective components in a batch of $n$
\item
  Number of successful packet transmissions out of $n$ attempts
\item
  Number of sensors that detect a signal (out of $n$ sensors)
\end{itemize}

\begin{PSbox}[unbreakable]{def}{Mathematical Definition}

$X \sim \mathrm{Binomial}(n,\, p)$

\end{PSbox}

\PSsubsection{PMF}

\PSdisplay{p_X(k) = \binom{n}{k} p^k (1-p)^{n-k}, \quad k = 0, 1, 2, \ldots, n}

where $C(n,\,k) = \frac{n!}{k!(n-k)!}$

\PSsubsection{CDF}

\PSdisplay{F_X(k) = \sum_{i=0}^{\lfloor k \rfloor} \binom{n}{i} p^i (1-p)^{n-i}}

No closed-form expression — use tables or MATLAB.

\PSsubsection{Mean}

\PSdisplay{E[X] = np}

\PSsubsection{Variance}

\PSdisplay{\text{Var}(X) = np(1-p)}

\PSsubsection{When to Use}

\PSyes{} Fixed number of trials $n$\PSbr{}\PSyes{} Each trial is independent\PSbr{}\PSyes{} Each trial has the same probability $p$\PSbr{}\PSyes{} Counting the number of ``successes''

\PSsubsection{When NOT to Use}

\PSno{} Trials are not independent (correlated failures)\PSbr{}\PSno{} Probability changes between trials\PSbr{}\PSno{} $n$ is not fixed (use Negative Binomial or Poisson)\PSbr{}\PSno{} Very large $n$ with small $p$ (use Poisson approximation)

\PSsubsection{MATLAB Implementation}

\tcbset{PScodemode/.style={unbreakable}}

\begin{Shaded}
\begin{Highlighting}[]
\CommentTok{\%\% Binomial Distribution: Bit Errors in a Frame}
\VariableTok{n} \OperatorTok{=} \FloatTok{100}\OperatorTok{;}         \CommentTok{\% bits per frame}
\VariableTok{p} \OperatorTok{=} \FloatTok{0.05}\OperatorTok{;}        \CommentTok{\% bit error probability}
\VariableTok{k} \OperatorTok{=} \FloatTok{0}\OperatorTok{:}\FloatTok{20}\OperatorTok{;}        \CommentTok{\% range of interest}

\CommentTok{\% PMF and CDF}
\VariableTok{pmf} \OperatorTok{=} \VariableTok{binopdf}\NormalTok{(}\VariableTok{k}\OperatorTok{,} \VariableTok{n}\OperatorTok{,} \VariableTok{p}\NormalTok{)}\OperatorTok{;}
\VariableTok{cdf\_vals} \OperatorTok{=} \VariableTok{binocdf}\NormalTok{(}\VariableTok{k}\OperatorTok{,} \VariableTok{n}\OperatorTok{,} \VariableTok{p}\NormalTok{)}\OperatorTok{;}

\CommentTok{\% Plot PMF}
\VariableTok{figure}\OperatorTok{;}
\VariableTok{stem}\NormalTok{(}\VariableTok{k}\OperatorTok{,} \VariableTok{pmf}\OperatorTok{,} \SpecialStringTok{\textquotesingle{}filled\textquotesingle{}}\OperatorTok{,} \SpecialStringTok{\textquotesingle{}LineWidth\textquotesingle{}}\OperatorTok{,} \FloatTok{1.5}\NormalTok{)}\OperatorTok{;}
\VariableTok{xlabel}\NormalTok{(}\SpecialStringTok{\textquotesingle{}Number of Errors (k)\textquotesingle{}}\NormalTok{)}\OperatorTok{;}
\VariableTok{ylabel}\NormalTok{(}\SpecialStringTok{\textquotesingle{}P(X = k)\textquotesingle{}}\NormalTok{)}\OperatorTok{;}
\VariableTok{title}\NormalTok{(}\VariableTok{sprintf}\NormalTok{(}\SpecialStringTok{\textquotesingle{}Binomial PMF: n=\%d, p=\%.2f\textquotesingle{}}\OperatorTok{,} \VariableTok{n}\OperatorTok{,} \VariableTok{p}\NormalTok{))}\OperatorTok{;}
\VariableTok{grid} \VariableTok{on}\OperatorTok{;}

\CommentTok{\% Simulation verification}
\VariableTok{N\_sim} \OperatorTok{=} \FloatTok{100000}\OperatorTok{;}
\VariableTok{X\_sim} \OperatorTok{=} \VariableTok{binornd}\NormalTok{(}\VariableTok{n}\OperatorTok{,} \VariableTok{p}\OperatorTok{,} \FloatTok{1}\OperatorTok{,} \VariableTok{N\_sim}\NormalTok{)}\OperatorTok{;}
\VariableTok{fprintf}\NormalTok{(}\SpecialStringTok{\textquotesingle{}E[X]: Theory = \%.2f, Simulation = \%.2f\textbackslash{}n\textquotesingle{}}\OperatorTok{,} \VariableTok{n}\OperatorTok{*}\VariableTok{p}\OperatorTok{,} \VariableTok{mean}\NormalTok{(}\VariableTok{X\_sim}\NormalTok{))}\OperatorTok{;}
\VariableTok{fprintf}\NormalTok{(}\SpecialStringTok{\textquotesingle{}Var(X): Theory = \%.2f, Simulation = \%.2f\textbackslash{}n\textquotesingle{}}\OperatorTok{,} \VariableTok{n}\OperatorTok{*}\VariableTok{p}\OperatorTok{*}\NormalTok{(}\FloatTok{1}\OperatorTok{{-}}\VariableTok{p}\NormalTok{)}\OperatorTok{,} \VariableTok{var}\NormalTok{(}\VariableTok{X\_sim}\NormalTok{))}\OperatorTok{;}
\end{Highlighting}
\end{Shaded}

\PSsubsection{Simulation Example: Packet Error Rate}

\tcbset{PScodemode/.style={unbreakable}}

\begin{Shaded}
\begin{Highlighting}[]
\CommentTok{\%\% Engineering Scenario: How many packets out of 50 have errors?}
\CommentTok{\% Each packet has 1000 bits, BER = 1e{-}4}
\CommentTok{\% A packet has an error if ANY bit is in error.}

\VariableTok{BER} \OperatorTok{=} \FloatTok{1e{-}4}\OperatorTok{;}
\VariableTok{bits\_per\_packet} \OperatorTok{=} \FloatTok{1000}\OperatorTok{;}
\VariableTok{n\_packets} \OperatorTok{=} \FloatTok{50}\OperatorTok{;}

\CommentTok{\% P(packet error) = 1 {-} P(all bits correct)}
\VariableTok{p\_packet\_error} \OperatorTok{=} \FloatTok{1} \OperatorTok{{-}}\NormalTok{ (}\FloatTok{1} \OperatorTok{{-}} \VariableTok{BER}\NormalTok{)}\OperatorTok{\^{}}\VariableTok{bits\_per\_packet}\OperatorTok{;}
\VariableTok{fprintf}\NormalTok{(}\SpecialStringTok{\textquotesingle{}P(packet error) = \%.4f\textbackslash{}n\textquotesingle{}}\OperatorTok{,} \VariableTok{p\_packet\_error}\NormalTok{)}\OperatorTok{;}

\CommentTok{\% Number of errored packets \textasciitilde{} Binomial(50, p\_packet\_error)}
\VariableTok{k} \OperatorTok{=} \FloatTok{0}\OperatorTok{:}\FloatTok{10}\OperatorTok{;}
\VariableTok{pmf\_packets} \OperatorTok{=} \VariableTok{binopdf}\NormalTok{(}\VariableTok{k}\OperatorTok{,} \VariableTok{n\_packets}\OperatorTok{,} \VariableTok{p\_packet\_error}\NormalTok{)}\OperatorTok{;}
\VariableTok{fprintf}\NormalTok{(}\SpecialStringTok{\textquotesingle{}P(0 packet errors) = \%.4f\textbackslash{}n\textquotesingle{}}\OperatorTok{,} \VariableTok{pmf\_packets}\NormalTok{(}\FloatTok{1}\NormalTok{))}\OperatorTok{;}
\VariableTok{fprintf}\NormalTok{(}\SpecialStringTok{\textquotesingle{}P(≤ 2 packet errors) = \%.4f\textbackslash{}n\textquotesingle{}}\OperatorTok{,} \VariableTok{sum}\NormalTok{(}\VariableTok{pmf\_packets}\NormalTok{(}\FloatTok{1}\OperatorTok{:}\FloatTok{3}\NormalTok{)))}\OperatorTok{;}
\end{Highlighting}
\end{Shaded}

\PSsection{3.4}{Geometric Distribution}

\PSsubsection{Physical Interpretation}

Models the \textbf{number of trials until the first success} (or equivalently, the number of failures before the first success).

\PSsubsection{Engineering Context}

\begin{itemize}
\tightlist
\item
  Number of transmissions until a packet is successfully delivered
\item
  Number of components tested until finding a defective one
\item
  Number of login attempts until successful authentication
\item
  Number of time slots until channel becomes available
\end{itemize}

\begin{PSbox}[unbreakable]{def}{Mathematical Definition}

$X \sim \mathrm{Geometric}(p)$

Convention: $X$ = number of the trial on which first success occurs $(X \in \{1,\, 2,\, 3,\, \ldots \})$.

\end{PSbox}

\PSsubsection{PMF}

\PSdisplay{p_X(k) = (1-p)^{k-1} p, \quad k = 1, 2, 3, \ldots}

($k-1$ failures, then one success)

\PSsubsection{CDF}

\PSdisplay{F_X(k) = 1 - (1-p)^k, \quad k = 1, 2, 3, \ldots}

\PSsubsection{Mean}

\PSdisplay{E[X] = \frac{1}{p}}

\PSsubsection{Variance}

\PSdisplay{\text{Var}(X) = \frac{1-p}{p^2}}

\PSsubsection{Memoryless Property}

The geometric distribution is the ONLY discrete distribution with the memoryless property:

\PSdisplay{P(X > m + n \mid X > m) = P(X > n)}

\textbf{Engineering meaning:} If a system hasn’t succeeded in $m$ attempts, the probability of needing $n$ more attempts is the same as starting fresh. ``Past failures don’t affect future probability.''

\PSsubsection{When to Use}

\PSyes{} Repeated independent trials until first success\PSbr{}\PSyes{} Each trial has the same probability\PSbr{}\PSyes{} Memoryless assumption is appropriate (e.g., independent retransmissions)

\PSsubsection{When NOT to Use}

\PSno{} Success probability changes over time (aging, learning)\PSbr{}\PSno{} Trials are not independent\PSbr{}\PSno{} Looking for the $r-\mathrm{th}$ success (use Negative Binomial)

\PSsubsection{MATLAB Implementation}

\tcbset{PScodemode/.style={unbreakable}}

\begin{Shaded}
\begin{Highlighting}[]
\CommentTok{\%\% Geometric Distribution: Retransmissions Until Success}
\VariableTok{p\_success} \OperatorTok{=} \FloatTok{0.7}\OperatorTok{;}   \CommentTok{\% P(successful transmission)}
\VariableTok{k} \OperatorTok{=} \FloatTok{1}\OperatorTok{:}\FloatTok{15}\OperatorTok{;}

\CommentTok{\% PMF}
\VariableTok{pmf\_geom} \OperatorTok{=} \VariableTok{geopdf}\NormalTok{(}\VariableTok{k}\OperatorTok{{-}}\FloatTok{1}\OperatorTok{,} \VariableTok{p\_success}\NormalTok{)}\OperatorTok{;}  \CommentTok{\% MATLAB counts failures (k{-}1)}

\CommentTok{\% Simulation}
\VariableTok{N} \OperatorTok{=} \FloatTok{100000}\OperatorTok{;}
\VariableTok{X\_sim} \OperatorTok{=} \VariableTok{geornd}\NormalTok{(}\VariableTok{p\_success}\OperatorTok{,} \FloatTok{1}\OperatorTok{,} \VariableTok{N}\NormalTok{) }\OperatorTok{+} \FloatTok{1}\OperatorTok{;}  \CommentTok{\% +1 to count trial number}

\VariableTok{figure}\OperatorTok{;}
\VariableTok{subplot}\NormalTok{(}\FloatTok{1}\OperatorTok{,}\FloatTok{2}\OperatorTok{,}\FloatTok{1}\NormalTok{)}\OperatorTok{;}
\VariableTok{stem}\NormalTok{(}\VariableTok{k}\OperatorTok{,} \VariableTok{pmf\_geom}\OperatorTok{,} \SpecialStringTok{\textquotesingle{}filled\textquotesingle{}}\NormalTok{)}\OperatorTok{;}
\VariableTok{xlabel}\NormalTok{(}\SpecialStringTok{\textquotesingle{}Trial Number k\textquotesingle{}}\NormalTok{)}\OperatorTok{;} \VariableTok{ylabel}\NormalTok{(}\SpecialStringTok{\textquotesingle{}P(X = k)\textquotesingle{}}\NormalTok{)}\OperatorTok{;}
\VariableTok{title}\NormalTok{(}\VariableTok{sprintf}\NormalTok{(}\SpecialStringTok{\textquotesingle{}Geometric PMF (p = \%.1f)\textquotesingle{}}\OperatorTok{,} \VariableTok{p\_success}\NormalTok{))}\OperatorTok{;}
\VariableTok{grid} \VariableTok{on}\OperatorTok{;}

\VariableTok{subplot}\NormalTok{(}\FloatTok{1}\OperatorTok{,}\FloatTok{2}\OperatorTok{,}\FloatTok{2}\NormalTok{)}\OperatorTok{;}
\VariableTok{histogram}\NormalTok{(}\VariableTok{X\_sim}\OperatorTok{,} \SpecialStringTok{\textquotesingle{}Normalization\textquotesingle{}}\OperatorTok{,} \SpecialStringTok{\textquotesingle{}probability\textquotesingle{}}\OperatorTok{,} \SpecialStringTok{\textquotesingle{}BinMethod\textquotesingle{}}\OperatorTok{,} \SpecialStringTok{\textquotesingle{}integers\textquotesingle{}}\NormalTok{)}\OperatorTok{;}
\VariableTok{xlabel}\NormalTok{(}\SpecialStringTok{\textquotesingle{}Trial Number k\textquotesingle{}}\NormalTok{)}\OperatorTok{;} \VariableTok{ylabel}\NormalTok{(}\SpecialStringTok{\textquotesingle{}Relative Frequency\textquotesingle{}}\NormalTok{)}\OperatorTok{;}
\VariableTok{title}\NormalTok{(}\SpecialStringTok{\textquotesingle{}Simulation Histogram\textquotesingle{}}\NormalTok{)}\OperatorTok{;}
\VariableTok{grid} \VariableTok{on}\OperatorTok{;}

\VariableTok{fprintf}\NormalTok{(}\SpecialStringTok{\textquotesingle{}E[X]: Theory = \%.2f, Simulation = \%.2f\textbackslash{}n\textquotesingle{}}\OperatorTok{,} \FloatTok{1}\OperatorTok{/}\VariableTok{p\_success}\OperatorTok{,} \VariableTok{mean}\NormalTok{(}\VariableTok{X\_sim}\NormalTok{))}\OperatorTok{;}
\end{Highlighting}
\end{Shaded}

\PSsection{3.5}{Negative Binomial Distribution}

\PSsubsection{Physical Interpretation}

Models the \textbf{number of trials until the $r-\mathrm{th}$ success} (generalization of Geometric).

\PSsubsection{Engineering Context}

\begin{itemize}
\tightlist
\item
  Number of transmissions until $r$ packets are successfully delivered
\item
  Number of components inspected until finding $r$ defectives
\item
  Number of time slots until $r$ channels become available
\end{itemize}

\begin{PSbox}[unbreakable]{def}{Mathematical Definition}

$X \sim \mathrm{NegBin}(r,\, p)$ = number of trials until $r-\mathrm{th}$ success.

\end{PSbox}

\PSsubsection{PMF}

\PSdisplay{p_X(k) = \binom{k-1}{r-1} p^r (1-p)^{k-r}, \quad k = r, r+1, r+2, \ldots}

(Choose which $r-1$ of the first $k-1$ trials were successes, the $k-\mathrm{th}$ is the $r-\mathrm{th}$ success.)

\PSsubsection{CDF}

No simple closed form. Computed by summation or MATLAB.

\PSsubsection{Mean}

\PSdisplay{E[X] = \frac{r}{p}}

\PSsubsection{Variance}

\PSdisplay{\text{Var}(X) = \frac{r(1-p)}{p^2}}

\PSsubsection{Relationship to Geometric}

When $r = 1$, the Negative Binomial reduces to the Geometric distribution.

\PSsubsection{When to Use}

\PSyes{} Repeated independent trials until $r-\mathrm{th}$ success\PSbr{}\PSyes{} Constant probability per trial\PSbr{}\PSyes{} Generalizing geometric to multiple required successes

\PSsubsection{When NOT to Use}

\PSno{} Only need first success (use Geometric — simpler)\PSbr{}\PSno{} Trials not independent\PSbr{}\PSno{} Fixed number of trials (use Binomial)

\PSsubsection{MATLAB Implementation}

\tcbset{PScodemode/.style={unbreakable}}

\begin{Shaded}
\begin{Highlighting}[]
\CommentTok{\%\% Negative Binomial: Transmissions Until r Successes}
\VariableTok{r} \OperatorTok{=} \FloatTok{5}\OperatorTok{;}             \CommentTok{\% Need 5 successful deliveries}
\VariableTok{p\_success} \OperatorTok{=} \FloatTok{0.8}\OperatorTok{;}   \CommentTok{\% P(success per trial)}

\CommentTok{\% MATLAB\textquotesingle{}s nbinpdf counts FAILURES before r{-}th success}
\CommentTok{\% X\_matlab = number of failures, so total trials = X\_matlab + r}
\VariableTok{k\_failures} \OperatorTok{=} \FloatTok{0}\OperatorTok{:}\FloatTok{20}\OperatorTok{;}

\VariableTok{pmf\_nb} \OperatorTok{=} \VariableTok{nbinpdf}\NormalTok{(}\VariableTok{k\_failures}\OperatorTok{,} \VariableTok{r}\OperatorTok{,} \VariableTok{p\_success}\NormalTok{)}\OperatorTok{;}

\VariableTok{figure}\OperatorTok{;}
\VariableTok{stem}\NormalTok{(}\VariableTok{k\_failures} \OperatorTok{+} \VariableTok{r}\OperatorTok{,} \VariableTok{pmf\_nb}\OperatorTok{,} \SpecialStringTok{\textquotesingle{}filled\textquotesingle{}}\NormalTok{)}\OperatorTok{;}
\VariableTok{xlabel}\NormalTok{(}\SpecialStringTok{\textquotesingle{}Total Trials Until 5th Success\textquotesingle{}}\NormalTok{)}\OperatorTok{;}
\VariableTok{ylabel}\NormalTok{(}\SpecialStringTok{\textquotesingle{}Probability\textquotesingle{}}\NormalTok{)}\OperatorTok{;}
\VariableTok{title}\NormalTok{(}\VariableTok{sprintf}\NormalTok{(}\SpecialStringTok{\textquotesingle{}Negative Binomial: r=\%d, p=\%.1f\textquotesingle{}}\OperatorTok{,} \VariableTok{r}\OperatorTok{,} \VariableTok{p\_success}\NormalTok{))}\OperatorTok{;}
\VariableTok{grid} \VariableTok{on}\OperatorTok{;}

\CommentTok{\% Simulation}
\VariableTok{N} \OperatorTok{=} \FloatTok{100000}\OperatorTok{;}
\VariableTok{X\_sim} \OperatorTok{=} \VariableTok{nbinrnd}\NormalTok{(}\VariableTok{r}\OperatorTok{,} \VariableTok{p\_success}\OperatorTok{,} \FloatTok{1}\OperatorTok{,} \VariableTok{N}\NormalTok{) }\OperatorTok{+} \VariableTok{r}\OperatorTok{;}  \CommentTok{\% total trials}
\VariableTok{fprintf}\NormalTok{(}\SpecialStringTok{\textquotesingle{}E[trials]: Theory = \%.2f, Simulation = \%.2f\textbackslash{}n\textquotesingle{}}\OperatorTok{,} \VariableTok{r}\OperatorTok{/}\VariableTok{p\_success}\OperatorTok{,} \VariableTok{mean}\NormalTok{(}\VariableTok{X\_sim}\NormalTok{))}\OperatorTok{;}
\end{Highlighting}
\end{Shaded}

\PSsection{3.6}{Poisson Distribution}

\PSsubsection{Physical Interpretation}

Models the \textbf{number of events occurring in a fixed interval} (time, space, area) when events occur independently at a constant average rate.

\PSsubsection{Engineering Context}

\begin{itemize}
\tightlist
\item
  Number of photons hitting a detector in $1 \mu s$
\item
  Number of packets arriving at a router in 1 second
\item
  Number of defects per $cm^{2}$ on a silicon wafer
\item
  Number of cosmic ray events in a sensor per hour
\item
  Number of interference spikes in a measurement window
\end{itemize}

\begin{PSbox}[unbreakable]{def}{Mathematical Definition}

$X \sim \mathrm{Poisson}(\lambda)$, where $\lambda$ = average number of events per interval.

\end{PSbox}

\PSsubsection{PMF}

\PSdisplay{p_X(k) = \frac{\lambda^k e^{-\lambda}}{k!}, \quad k = 0, 1, 2, \ldots}

\PSsubsection{CDF}

\PSdisplay{F_X(k) = e^{-\lambda} \sum_{i=0}^{\lfloor k \rfloor} \frac{\lambda^i}{i!}}

\PSsubsection{Mean}

\PSdisplay{E[X] = \lambda}

\PSsubsection{Variance}

\PSdisplay{\text{Var}(X) = \lambda}

\textbf{Unique property:} Mean equals variance! This is a quick check — if sample mean \ensuremath{\approx} sample variance, data may be Poisson.

\PSsubsection{Poisson as Binomial Approximation}

When $n$ is large and $p$ is small, with $\lambda$ = np:

\PSdisplay{\text{Binomial}(n, p) \approx \text{Poisson}(\lambda = np)}

Rule of thumb: $n \ge 20$ and $p \le 0.05$.

\PSsubsection{When to Use}

\PSyes{} Counting events in a fixed interval\PSbr{}\PSyes{} Events occur independently\PSbr{}\PSyes{} Events occur at a constant average rate\PSbr{}\PSyes{} Rare events with many opportunities (large $n$, small $p$)\PSbr{}\PSyes{} Mean \ensuremath{\approx} Variance in data

\PSsubsection{When NOT to Use}

\PSno{} Events are not independent (clustered arrivals)\PSbr{}\PSno{} Rate varies with time (non-homogeneous process)\PSbr{}\PSno{} Events are not rare ($p$ not small) — use Binomial\PSbr{}\PSno{} Mean \ensuremath{\neq} Variance (look at overdispersion/underdispersion)

\PSsubsection{MATLAB Implementation}

\tcbset{PScodemode/.style={unbreakable}}

\begin{Shaded}
\begin{Highlighting}[]
\CommentTok{\%\% Poisson Distribution: Network Packet Arrivals}
\VariableTok{lambda} \OperatorTok{=} \FloatTok{4.5}\OperatorTok{;}     \CommentTok{\% Average packets per millisecond}
\VariableTok{k} \OperatorTok{=} \FloatTok{0}\OperatorTok{:}\FloatTok{15}\OperatorTok{;}

\VariableTok{pmf\_pois} \OperatorTok{=} \VariableTok{poisspdf}\NormalTok{(}\VariableTok{k}\OperatorTok{,} \VariableTok{lambda}\NormalTok{)}\OperatorTok{;}
\VariableTok{cdf\_pois} \OperatorTok{=} \VariableTok{poisscdf}\NormalTok{(}\VariableTok{k}\OperatorTok{,} \VariableTok{lambda}\NormalTok{)}\OperatorTok{;}

\VariableTok{figure}\OperatorTok{;}
\VariableTok{subplot}\NormalTok{(}\FloatTok{1}\OperatorTok{,}\FloatTok{2}\OperatorTok{,}\FloatTok{1}\NormalTok{)}\OperatorTok{;}
\VariableTok{stem}\NormalTok{(}\VariableTok{k}\OperatorTok{,} \VariableTok{pmf\_pois}\OperatorTok{,} \SpecialStringTok{\textquotesingle{}filled\textquotesingle{}}\OperatorTok{,} \SpecialStringTok{\textquotesingle{}LineWidth\textquotesingle{}}\OperatorTok{,} \FloatTok{1.5}\NormalTok{)}\OperatorTok{;}
\VariableTok{xlabel}\NormalTok{(}\SpecialStringTok{\textquotesingle{}Number of Arrivals (k)\textquotesingle{}}\NormalTok{)}\OperatorTok{;}
\VariableTok{ylabel}\NormalTok{(}\SpecialStringTok{\textquotesingle{}P(X = k)\textquotesingle{}}\NormalTok{)}\OperatorTok{;}
\VariableTok{title}\NormalTok{(}\VariableTok{sprintf}\NormalTok{(}\SpecialStringTok{\textquotesingle{}Poisson PMF (λ = \%.1f)\textquotesingle{}}\OperatorTok{,} \VariableTok{lambda}\NormalTok{))}\OperatorTok{;}
\VariableTok{grid} \VariableTok{on}\OperatorTok{;}

\VariableTok{subplot}\NormalTok{(}\FloatTok{1}\OperatorTok{,}\FloatTok{2}\OperatorTok{,}\FloatTok{2}\NormalTok{)}\OperatorTok{;}
\VariableTok{stairs}\NormalTok{(}\VariableTok{k}\OperatorTok{,} \VariableTok{cdf\_pois}\OperatorTok{,} \SpecialStringTok{\textquotesingle{}LineWidth\textquotesingle{}}\OperatorTok{,} \FloatTok{2}\NormalTok{)}\OperatorTok{;}
\VariableTok{xlabel}\NormalTok{(}\SpecialStringTok{\textquotesingle{}k\textquotesingle{}}\NormalTok{)}\OperatorTok{;} \VariableTok{ylabel}\NormalTok{(}\SpecialStringTok{\textquotesingle{}P(X ≤ k)\textquotesingle{}}\NormalTok{)}\OperatorTok{;}
\VariableTok{title}\NormalTok{(}\SpecialStringTok{\textquotesingle{}Poisson CDF\textquotesingle{}}\NormalTok{)}\OperatorTok{;}
\VariableTok{grid} \VariableTok{on}\OperatorTok{;}
\end{Highlighting}
\end{Shaded}

\PSsubsection{Simulation Example: Photon Counting}

\tcbset{PScodemode/.style={unbreakable}}

\begin{Shaded}
\begin{Highlighting}[]
\CommentTok{\%\% Engineering Scenario: Photon Detection}
\CommentTok{\% An optical receiver detects photons. Average rate: 10 photons/μs.}
\CommentTok{\% What is P(detecting fewer than 5 photons in one μs)?}

\VariableTok{lambda} \OperatorTok{=} \FloatTok{10}\OperatorTok{;}   \CommentTok{\% average photons per μs}
\VariableTok{N} \OperatorTok{=} \FloatTok{100000}\OperatorTok{;}

\CommentTok{\% Analytical}
\VariableTok{p\_fewer\_5} \OperatorTok{=} \VariableTok{poisscdf}\NormalTok{(}\FloatTok{4}\OperatorTok{,} \VariableTok{lambda}\NormalTok{)}\OperatorTok{;}
\VariableTok{fprintf}\NormalTok{(}\SpecialStringTok{\textquotesingle{}P(X \textless{} 5): Analytical = \%.6f\textbackslash{}n\textquotesingle{}}\OperatorTok{,} \VariableTok{p\_fewer\_5}\NormalTok{)}\OperatorTok{;}

\CommentTok{\% Simulation}
\VariableTok{X\_sim} \OperatorTok{=} \VariableTok{poissrnd}\NormalTok{(}\VariableTok{lambda}\OperatorTok{,} \FloatTok{1}\OperatorTok{,} \VariableTok{N}\NormalTok{)}\OperatorTok{;}
\VariableTok{p\_fewer\_5\_sim} \OperatorTok{=} \VariableTok{mean}\NormalTok{(}\VariableTok{X\_sim} \OperatorTok{\textless{}} \FloatTok{5}\NormalTok{)}\OperatorTok{;}
\VariableTok{fprintf}\NormalTok{(}\SpecialStringTok{\textquotesingle{}P(X \textless{} 5): Simulation = \%.6f\textbackslash{}n\textquotesingle{}}\OperatorTok{,} \VariableTok{p\_fewer\_5\_sim}\NormalTok{)}\OperatorTok{;}

\CommentTok{\% Verify mean = variance (Poisson signature)}
\VariableTok{fprintf}\NormalTok{(}\SpecialStringTok{\textquotesingle{}Sample Mean = \%.3f, Sample Variance = \%.3f\textbackslash{}n\textquotesingle{}}\OperatorTok{,} \VariableTok{mean}\NormalTok{(}\VariableTok{X\_sim}\NormalTok{)}\OperatorTok{,} \VariableTok{var}\NormalTok{(}\VariableTok{X\_sim}\NormalTok{))}\OperatorTok{;}
\end{Highlighting}
\end{Shaded}

\PSsection{3.7}{Uniform Distribution (Continuous)}

\PSsubsection{Physical Interpretation}

Models a quantity that is \textbf{equally likely} to take any value in an interval $[a,\, b]$. Represents \textbf{maximum uncertainty} within a known range.

\PSsubsection{Engineering Context}

\begin{itemize}
\tightlist
\item
  Phase of a received signal (unknown, equally likely in $[0,\, 2\pi)$)
\item
  Round-off error in an ADC (uniform in $\left[-\frac{\Delta}{2},\, \frac{\Delta}{2}\right]$ where $\Delta$ = quantization step)
\item
  Arrival time within a time slot (unknown position)
\item
  Random access protocol: choose a random backoff time
\end{itemize}

\begin{PSbox}[unbreakable]{def}{Mathematical Definition}

$X \sim \mathrm{Uniform}(a,\, b)$ or $X \sim U(a,\, b)$

\end{PSbox}

\PSsubsection{PDF}

\PSdisplay{f_X(x) = \begin{cases} \frac{1}{b-a} & a \leq x \leq b \\ 0 & \text{otherwise} \end{cases}}

\PSsubsection{CDF}

\PSdisplay{F_X(x) = \begin{cases} 0 & x < a \\ \frac{x-a}{b-a} & a \leq x \leq b \\ 1 & x > b \end{cases}}

\PSsubsection{Mean}

\PSdisplay{E[X] = \frac{a+b}{2}}

\PSsubsection{Variance}

\PSdisplay{\text{Var}(X) = \frac{(b-a)^2}{12}}

\PSsubsection{When to Use}

\PSyes{} All values in a range are equally likely\PSbr{}\PSyes{} No information favors any particular value\PSbr{}\PSyes{} Quantization error modeling\PSbr{}\PSyes{} Phase of unmodulated carrier

\PSsubsection{When NOT to Use}

\PSno{} Values near the center are more likely (use Gaussian)\PSbr{}\PSno{} Values are non-negative with exponential decay (use Exponential)\PSbr{}\PSno{} Range is unbounded

\PSsubsection{MATLAB Implementation}

\tcbset{PScodemode/.style={unbreakable}}

\begin{Shaded}
\begin{Highlighting}[]
\CommentTok{\%\% Uniform Distribution: Quantization Error}
\VariableTok{a} \OperatorTok{=} \OperatorTok{{-}}\FloatTok{0.5}\OperatorTok{;}  \VariableTok{b} \OperatorTok{=} \FloatTok{0.5}\OperatorTok{;}  \CommentTok{\% Quantization step = 1 LSB}
\VariableTok{N} \OperatorTok{=} \FloatTok{100000}\OperatorTok{;}

\VariableTok{X} \OperatorTok{=} \VariableTok{a} \OperatorTok{+}\NormalTok{ (}\VariableTok{b}\OperatorTok{{-}}\VariableTok{a}\NormalTok{)}\OperatorTok{*}\VariableTok{rand}\NormalTok{(}\FloatTok{1}\OperatorTok{,} \VariableTok{N}\NormalTok{)}\OperatorTok{;}  \CommentTok{\% or: unifrnd(a, b, 1, N)}

\VariableTok{figure}\OperatorTok{;}
\VariableTok{histogram}\NormalTok{(}\VariableTok{X}\OperatorTok{,} \FloatTok{50}\OperatorTok{,} \SpecialStringTok{\textquotesingle{}Normalization\textquotesingle{}}\OperatorTok{,} \SpecialStringTok{\textquotesingle{}pdf\textquotesingle{}}\NormalTok{)}\OperatorTok{;}
\VariableTok{hold} \VariableTok{on}\OperatorTok{;}
\VariableTok{plot}\NormalTok{([}\VariableTok{a} \VariableTok{a} \VariableTok{b} \VariableTok{b}\NormalTok{]}\OperatorTok{,}\NormalTok{ [}\FloatTok{0} \FloatTok{1}\OperatorTok{/}\NormalTok{(}\VariableTok{b}\OperatorTok{{-}}\VariableTok{a}\NormalTok{) }\FloatTok{1}\OperatorTok{/}\NormalTok{(}\VariableTok{b}\OperatorTok{{-}}\VariableTok{a}\NormalTok{) }\FloatTok{0}\NormalTok{]}\OperatorTok{,} \SpecialStringTok{\textquotesingle{}r{-}\textquotesingle{}}\OperatorTok{,} \SpecialStringTok{\textquotesingle{}LineWidth\textquotesingle{}}\OperatorTok{,} \FloatTok{2}\NormalTok{)}\OperatorTok{;}
\VariableTok{xlabel}\NormalTok{(}\SpecialStringTok{\textquotesingle{}Quantization Error\textquotesingle{}}\NormalTok{)}\OperatorTok{;} \VariableTok{ylabel}\NormalTok{(}\SpecialStringTok{\textquotesingle{}f\_X(x)\textquotesingle{}}\NormalTok{)}\OperatorTok{;}
\VariableTok{title}\NormalTok{(}\SpecialStringTok{\textquotesingle{}Uniform PDF: Quantization Error\textquotesingle{}}\NormalTok{)}\OperatorTok{;}
\VariableTok{legend}\NormalTok{(}\SpecialStringTok{\textquotesingle{}Simulated\textquotesingle{}}\OperatorTok{,} \SpecialStringTok{\textquotesingle{}Theoretical\textquotesingle{}}\NormalTok{)}\OperatorTok{;}

\VariableTok{fprintf}\NormalTok{(}\SpecialStringTok{\textquotesingle{}Mean: Theory=\%.3f, Sim=\%.3f\textbackslash{}n\textquotesingle{}}\OperatorTok{,}\NormalTok{ (}\VariableTok{a}\OperatorTok{+}\VariableTok{b}\NormalTok{)}\OperatorTok{/}\FloatTok{2}\OperatorTok{,} \VariableTok{mean}\NormalTok{(}\VariableTok{X}\NormalTok{))}\OperatorTok{;}
\VariableTok{fprintf}\NormalTok{(}\SpecialStringTok{\textquotesingle{}Var:  Theory=\%.4f, Sim=\%.4f\textbackslash{}n\textquotesingle{}}\OperatorTok{,}\NormalTok{ (}\VariableTok{b}\OperatorTok{{-}}\VariableTok{a}\NormalTok{)}\OperatorTok{\^{}}\FloatTok{2}\OperatorTok{/}\FloatTok{12}\OperatorTok{,} \VariableTok{var}\NormalTok{(}\VariableTok{X}\NormalTok{))}\OperatorTok{;}
\end{Highlighting}
\end{Shaded}

\PSsection{3.8}{Exponential Distribution}

\PSsubsection{Physical Interpretation}

Models the \textbf{time between events} in a Poisson process, or the \textbf{lifetime} of a memoryless component.

\PSsubsection{Engineering Context}

\begin{itemize}
\tightlist
\item
  Time between packet arrivals
\item
  Time until component failure (constant failure rate)
\item
  Time between cosmic ray events in a detector
\item
  Inter-arrival time of calls at a switch
\item
  Duration of a fading dip in a wireless channel
\end{itemize}

\begin{PSbox}[unbreakable]{def}{Mathematical Definition}

$X \sim \mathrm{Exponential}(\lambda)$ where $\lambda$ = rate parameter (events per unit time).

Alternative parameterization: $X \sim \mathrm{Exp}(\beta)$ where $\beta = \frac{1}{\lambda}$ = mean.

\end{PSbox}

\PSsubsection{PDF}

\PSdisplay{f_X(x) = \lambda e^{-\lambda x}, \quad x \geq 0}

Or with mean $\beta$: $f_{X}(x) = \frac{1}{\beta}e^{-x/\beta}$

\PSsubsection{CDF}

\PSdisplay{F_X(x) = 1 - e^{-\lambda x}, \quad x \geq 0}

\PSsubsection{Mean}

\PSdisplay{E[X] = \frac{1}{\lambda} = \beta}

\PSsubsection{Variance}

\PSdisplay{\text{Var}(X) = \frac{1}{\lambda^2} = \beta^2}

\textbf{Note:} Standard deviation = mean for exponential distribution.

\PSsubsection{Memoryless Property}

\PSdisplay{P(X > s + t \mid X > s) = P(X > t)}

The exponential is the ONLY continuous distribution with this property.

\textbf{Engineering meaning:} A component that has been working for $s$ hours is statistically ``as good as new.'' Its remaining lifetime has the same distribution as a brand-new component.

\PSsubsection{Relationship to Poisson}

If events arrive as a Poisson process with rate $\lambda$:

\begin{itemize}
\tightlist
\item
  Number of events in time $T \sim \mathrm{Poisson}(\lambda T)$
\item
  Time between events $\sim \mathrm{Exponential}(\lambda)$
\end{itemize}

\PSsubsection{When to Use}

\PSyes{} Modeling time between events in a Poisson process\PSbr{}\PSyes{} Component lifetime with constant failure rate\PSbr{}\PSyes{} Memoryless assumption is reasonable

\PSsubsection{When NOT to Use}

\PSno{} Failure rate increases with age (wear-out) — use Weibull or Gamma\PSbr{}\PSno{} Failure rate decreases with age (burn-in) — use Weibull\PSbr{}\PSno{} Data shows increasing hazard — NOT memoryless

\PSsubsection{MATLAB Implementation}

\tcbset{PScodemode/.style={unbreakable}}

\begin{Shaded}
\begin{Highlighting}[]
\CommentTok{\%\% Exponential Distribution: Component Lifetime}
\VariableTok{lambda} \OperatorTok{=} \FloatTok{0.001}\OperatorTok{;}          \CommentTok{\% Failure rate: 0.001 per hour}
\VariableTok{beta} \OperatorTok{=} \FloatTok{1}\OperatorTok{/}\VariableTok{lambda}\OperatorTok{;}         \CommentTok{\% Mean lifetime: 1000 hours}

\VariableTok{x} \OperatorTok{=} \VariableTok{linspace}\NormalTok{(}\FloatTok{0}\OperatorTok{,} \FloatTok{5000}\OperatorTok{,} \FloatTok{1000}\NormalTok{)}\OperatorTok{;}
\VariableTok{pdf\_exp} \OperatorTok{=} \VariableTok{exppdf}\NormalTok{(}\VariableTok{x}\OperatorTok{,} \VariableTok{beta}\NormalTok{)}\OperatorTok{;}
\VariableTok{cdf\_exp} \OperatorTok{=} \VariableTok{expcdf}\NormalTok{(}\VariableTok{x}\OperatorTok{,} \VariableTok{beta}\NormalTok{)}\OperatorTok{;}

\VariableTok{figure}\OperatorTok{;}
\VariableTok{subplot}\NormalTok{(}\FloatTok{2}\OperatorTok{,}\FloatTok{1}\OperatorTok{,}\FloatTok{1}\NormalTok{)}\OperatorTok{;}
\VariableTok{plot}\NormalTok{(}\VariableTok{x}\OperatorTok{,} \VariableTok{pdf\_exp}\OperatorTok{,} \SpecialStringTok{\textquotesingle{}b{-}\textquotesingle{}}\OperatorTok{,} \SpecialStringTok{\textquotesingle{}LineWidth\textquotesingle{}}\OperatorTok{,} \FloatTok{2}\NormalTok{)}\OperatorTok{;}
\VariableTok{xlabel}\NormalTok{(}\SpecialStringTok{\textquotesingle{}Time (hours)\textquotesingle{}}\NormalTok{)}\OperatorTok{;} \VariableTok{ylabel}\NormalTok{(}\SpecialStringTok{\textquotesingle{}f\_X(x)\textquotesingle{}}\NormalTok{)}\OperatorTok{;}
\VariableTok{title}\NormalTok{(}\VariableTok{sprintf}\NormalTok{(}\SpecialStringTok{\textquotesingle{}Exponential PDF: Mean Lifetime = \%d hours\textquotesingle{}}\OperatorTok{,} \VariableTok{beta}\NormalTok{))}\OperatorTok{;}
\VariableTok{grid} \VariableTok{on}\OperatorTok{;}

\VariableTok{subplot}\NormalTok{(}\FloatTok{2}\OperatorTok{,}\FloatTok{1}\OperatorTok{,}\FloatTok{2}\NormalTok{)}\OperatorTok{;}
\VariableTok{plot}\NormalTok{(}\VariableTok{x}\OperatorTok{,} \VariableTok{cdf\_exp}\OperatorTok{,} \SpecialStringTok{\textquotesingle{}r{-}\textquotesingle{}}\OperatorTok{,} \SpecialStringTok{\textquotesingle{}LineWidth\textquotesingle{}}\OperatorTok{,} \FloatTok{2}\NormalTok{)}\OperatorTok{;}
\VariableTok{xlabel}\NormalTok{(}\SpecialStringTok{\textquotesingle{}Time (hours)\textquotesingle{}}\NormalTok{)}\OperatorTok{;} \VariableTok{ylabel}\NormalTok{(}\SpecialStringTok{\textquotesingle{}F\_X(x) = P(failure by time x)\textquotesingle{}}\NormalTok{)}\OperatorTok{;}
\VariableTok{title}\NormalTok{(}\SpecialStringTok{\textquotesingle{}Exponential CDF\textquotesingle{}}\NormalTok{)}\OperatorTok{;}
\VariableTok{grid} \VariableTok{on}\OperatorTok{;}

\CommentTok{\% P(component lasts more than 2000 hours)}
\VariableTok{p\_survive\_2000} \OperatorTok{=} \FloatTok{1} \OperatorTok{{-}} \VariableTok{expcdf}\NormalTok{(}\FloatTok{2000}\OperatorTok{,} \VariableTok{beta}\NormalTok{)}\OperatorTok{;}
\VariableTok{fprintf}\NormalTok{(}\SpecialStringTok{\textquotesingle{}P(T \textgreater{} 2000 hours) = \%.4f\textbackslash{}n\textquotesingle{}}\OperatorTok{,} \VariableTok{p\_survive\_2000}\NormalTok{)}\OperatorTok{;}
\end{Highlighting}
\end{Shaded}

\PSsubsection{Simulation Example: Memoryless Property Verification}

\tcbset{PScodemode/.style={unbreakable}}

\begin{Shaded}
\begin{Highlighting}[]
\CommentTok{\%\% Verify Memoryless Property}
\VariableTok{beta} \OperatorTok{=} \FloatTok{100}\OperatorTok{;}   \CommentTok{\% Mean lifetime = 100 hours}
\VariableTok{N} \OperatorTok{=} \FloatTok{100000}\OperatorTok{;}

\VariableTok{X} \OperatorTok{=} \VariableTok{exprnd}\NormalTok{(}\VariableTok{beta}\OperatorTok{,} \FloatTok{1}\OperatorTok{,} \VariableTok{N}\NormalTok{)}\OperatorTok{;}

\CommentTok{\% Conditional: given X \textgreater{} 50, what is distribution of (X {-} 50)?}
\VariableTok{survivors} \OperatorTok{=} \VariableTok{X}\NormalTok{(}\VariableTok{X} \OperatorTok{\textgreater{}} \FloatTok{50}\NormalTok{)}\OperatorTok{;}
\VariableTok{remaining\_life} \OperatorTok{=} \VariableTok{survivors} \OperatorTok{{-}} \FloatTok{50}\OperatorTok{;}

\VariableTok{fprintf}\NormalTok{(}\SpecialStringTok{\textquotesingle{}Mean of X: \%.1f (should be \%.1f)\textbackslash{}n\textquotesingle{}}\OperatorTok{,} \VariableTok{mean}\NormalTok{(}\VariableTok{X}\NormalTok{)}\OperatorTok{,} \VariableTok{beta}\NormalTok{)}\OperatorTok{;}
\VariableTok{fprintf}\NormalTok{(}\SpecialStringTok{\textquotesingle{}Mean of remaining life given survival past 50: \%.1f\textbackslash{}n\textquotesingle{}}\OperatorTok{,} \VariableTok{mean}\NormalTok{(}\VariableTok{remaining\_life}\NormalTok{))}\OperatorTok{;}
\VariableTok{fprintf}\NormalTok{(}\SpecialStringTok{\textquotesingle{}Should also be \%.1f (memoryless!)\textbackslash{}n\textquotesingle{}}\OperatorTok{,} \VariableTok{beta}\NormalTok{)}\OperatorTok{;}
\end{Highlighting}
\end{Shaded}

\PSsection{3.9}{Gaussian (Normal) Distribution}

\PSsubsection{Physical Interpretation}

Models quantities that result from the \textbf{sum of many small, independent random effects}. The most important distribution in engineering due to the Central Limit Theorem.

\PSsubsection{Engineering Context}

\begin{itemize}
\tightlist
\item
  Thermal noise voltage (sum of many electron contributions)
\item
  Measurement errors (sum of many small error sources)
\item
  Manufacturing variations (many independent process variations)
\item
  Aggregate interference in communications
\item
  Any quantity well-modeled by CLT
\end{itemize}

\begin{PSbox}[unbreakable]{def}{Mathematical Definition}

$X \sim N(\mu,\, \sigma^{2})$ where $\mu$ = mean, $\sigma^{2}$ = variance.

\end{PSbox}

\PSsubsection{PDF}

\PSdisplay{f_X(x) = \frac{1}{\sigma\sqrt{2\pi}} \exp\left(-\frac{(x-\mu)^2}{2\sigma^2}\right), \quad -\infty < x < \infty}

\PSsubsection{CDF}

\PSdisplay{F_X(x) = \Phi\left(\frac{x-\mu}{\sigma}\right) = \frac{1}{2}\left[1 + \text{erf}\left(\frac{x-\mu}{\sigma\sqrt{2}}\right)\right]}

No closed-form — computed numerically or via tables.

\PSsubsection{Mean}

\PSdisplay{E[X] = \mu}

\PSsubsection{Variance}

\PSdisplay{\text{Var}(X) = \sigma^2}

\PSsubsection{Standard Normal}

$Z \sim N(0,\, 1)$: standardized form.

Any Normal can be standardized: $Z = \frac{X - \mu}{\sigma}$

\PSsubsection{Key Probability Rules $(\sigma -\text{rules})$}

{\def\LTcaptype{none} % do not increment counter
\begin{longtable}[c]{@{}cc@{}}
\toprule\noalign{}
\rowcolor{PSnavy}\PSth{Range} & \PSth{Probability} \\
\midrule\noalign{}
\endhead
\bottomrule\noalign{}
\endlastfoot
$\mu \pm 1\sigma$ & 68.27\% \\
$\mu \pm 2\sigma$ & 95.45\% \\
$\mu \pm 3\sigma$ & 99.73\% \\
$\mu \pm 4\sigma$ & 99.9937\% \\
\end{longtable}
}

\PSsubsection{Properties}

\begin{enumerate}
\def\labelenumi{\arabic{enumi}.}
\tightlist
\item
  \textbf{Symmetric} about $\mu$
\item
  \textbf{Linear transformation:} If $X \sim N(\mu,\, \sigma^{2})$, then aX $+ b \sim N(a\mu + b,\, a^{2}\sigma^{2})$
\item
  \textbf{Sum of Gaussians:} If $X \sim N(\mu_{1},\, \sigma_{1}^{2})$ and $Y \sim N(\mu_{2},\, \sigma_{2}^{2})$ are independent, then $X + Y \sim N(\mu_{1}+\mu_{2},\, \sigma_{1}^{2}+\sigma_{2}^{2})$
\item
  \textbf{Fully characterized} by mean and variance (all higher moments determined)
\end{enumerate}

\PSsubsection{When to Use}

\PSyes{} Thermal noise in circuits\PSbr{}\PSyes{} Measurement errors\PSbr{}\PSyes{} Any sum of many independent random effects (CLT)\PSbr{}\PSyes{} Manufacturing variations\PSbr{}\PSyes{} When data histogram is bell-shaped and symmetric

\PSsubsection{When NOT to Use}

\PSno{} Strictly positive quantities (Gaussian allows negative values)\PSbr{}\PSno{} Heavily skewed data\PSbr{}\PSno{} Bounded quantities (Gaussian has infinite support)\PSbr{}\PSno{} Discrete counts\PSbr{}\PSno{} Heavy-tailed phenomena (use t-distribution or Cauchy)

\PSsubsection{MATLAB Implementation}

\tcbset{PScodemode/.style={unbreakable}}

\begin{Shaded}
\begin{Highlighting}[]
\CommentTok{\%\% Gaussian Distribution: Thermal Noise}
\VariableTok{mu} \OperatorTok{=} \FloatTok{0}\OperatorTok{;}           \CommentTok{\% Zero{-}mean noise}
\VariableTok{sigma} \OperatorTok{=} \FloatTok{0.1}\OperatorTok{;}      \CommentTok{\% RMS noise = 100 mV}

\VariableTok{x} \OperatorTok{=} \VariableTok{linspace}\NormalTok{(}\VariableTok{mu} \OperatorTok{{-}} \FloatTok{4}\OperatorTok{*}\VariableTok{sigma}\OperatorTok{,} \VariableTok{mu} \OperatorTok{+} \FloatTok{4}\OperatorTok{*}\VariableTok{sigma}\OperatorTok{,} \FloatTok{1000}\NormalTok{)}\OperatorTok{;}
\VariableTok{pdf\_norm} \OperatorTok{=} \VariableTok{normpdf}\NormalTok{(}\VariableTok{x}\OperatorTok{,} \VariableTok{mu}\OperatorTok{,} \VariableTok{sigma}\NormalTok{)}\OperatorTok{;}

\VariableTok{figure}\OperatorTok{;}
\VariableTok{plot}\NormalTok{(}\VariableTok{x}\OperatorTok{,} \VariableTok{pdf\_norm}\OperatorTok{,} \SpecialStringTok{\textquotesingle{}b{-}\textquotesingle{}}\OperatorTok{,} \SpecialStringTok{\textquotesingle{}LineWidth\textquotesingle{}}\OperatorTok{,} \FloatTok{2}\NormalTok{)}\OperatorTok{;}
\VariableTok{hold} \VariableTok{on}\OperatorTok{;}
\CommentTok{\% Shade the ±2σ region}
\VariableTok{x\_fill} \OperatorTok{=} \VariableTok{linspace}\NormalTok{(}\VariableTok{mu}\OperatorTok{{-}}\FloatTok{2}\OperatorTok{*}\VariableTok{sigma}\OperatorTok{,} \VariableTok{mu}\OperatorTok{+}\FloatTok{2}\OperatorTok{*}\VariableTok{sigma}\OperatorTok{,} \FloatTok{200}\NormalTok{)}\OperatorTok{;}
\VariableTok{fill}\NormalTok{([}\VariableTok{x\_fill}\OperatorTok{,} \VariableTok{fliplr}\NormalTok{(}\VariableTok{x\_fill}\NormalTok{)]}\OperatorTok{,}\NormalTok{ [}\VariableTok{normpdf}\NormalTok{(}\VariableTok{x\_fill}\OperatorTok{,}\VariableTok{mu}\OperatorTok{,}\VariableTok{sigma}\NormalTok{)}\OperatorTok{,} \VariableTok{zeros}\NormalTok{(}\VariableTok{size}\NormalTok{(}\VariableTok{x\_fill}\NormalTok{))]}\OperatorTok{,} \OperatorTok{...}
     \SpecialStringTok{\textquotesingle{}b\textquotesingle{}}\OperatorTok{,} \SpecialStringTok{\textquotesingle{}FaceAlpha\textquotesingle{}}\OperatorTok{,} \FloatTok{0.3}\NormalTok{)}\OperatorTok{;}
\VariableTok{xlabel}\NormalTok{(}\SpecialStringTok{\textquotesingle{}Noise Voltage (V)\textquotesingle{}}\NormalTok{)}\OperatorTok{;}
\VariableTok{ylabel}\NormalTok{(}\SpecialStringTok{\textquotesingle{}f\_X(x)\textquotesingle{}}\NormalTok{)}\OperatorTok{;}
\VariableTok{title}\NormalTok{(}\VariableTok{sprintf}\NormalTok{(}\SpecialStringTok{\textquotesingle{}Gaussian PDF: μ=\%.1f, σ=\%.3f V\textquotesingle{}}\OperatorTok{,} \VariableTok{mu}\OperatorTok{,} \VariableTok{sigma}\NormalTok{))}\OperatorTok{;}
\VariableTok{legend}\NormalTok{(}\SpecialStringTok{\textquotesingle{}PDF\textquotesingle{}}\OperatorTok{,} \SpecialStringTok{\textquotesingle{}±2σ region (95.45\%)\textquotesingle{}}\NormalTok{)}\OperatorTok{;}
\VariableTok{grid} \VariableTok{on}\OperatorTok{;}

\CommentTok{\% Probability computations}
\VariableTok{fprintf}\NormalTok{(}\SpecialStringTok{\textquotesingle{}P(|noise| \textgreater{} 0.2V) = \%.4f\textbackslash{}n\textquotesingle{}}\OperatorTok{,} \FloatTok{2}\OperatorTok{*}\NormalTok{(}\FloatTok{1}\OperatorTok{{-}}\VariableTok{normcdf}\NormalTok{(}\FloatTok{0.2}\OperatorTok{,} \FloatTok{0}\OperatorTok{,} \VariableTok{sigma}\NormalTok{)))}\OperatorTok{;}
\VariableTok{fprintf}\NormalTok{(}\SpecialStringTok{\textquotesingle{}P(|noise| \textgreater{} 0.3V) = \%.4f\textbackslash{}n\textquotesingle{}}\OperatorTok{,} \FloatTok{2}\OperatorTok{*}\NormalTok{(}\FloatTok{1}\OperatorTok{{-}}\VariableTok{normcdf}\NormalTok{(}\FloatTok{0.3}\OperatorTok{,} \FloatTok{0}\OperatorTok{,} \VariableTok{sigma}\NormalTok{)))}\OperatorTok{;}
\end{Highlighting}
\end{Shaded}

\PSsubsection{Simulation Example: Q-Function and BER}

\tcbset{PScodemode/.style={unbreakable}}

\begin{Shaded}
\begin{Highlighting}[]
\CommentTok{\%\% Engineering Application: Binary Communication BER}
\CommentTok{\% In BPSK with AWGN: BER = Q(sqrt(2*Eb/N0))}
\CommentTok{\% Q(x) = 1 {-} Phi(x) = probability that N(0,1) \textgreater{} x}

\VariableTok{EbN0\_dB} \OperatorTok{=} \FloatTok{0}\OperatorTok{:}\FloatTok{0.5}\OperatorTok{:}\FloatTok{12}\OperatorTok{;}
\VariableTok{EbN0} \OperatorTok{=} \FloatTok{10}\OperatorTok{.\^{}}\NormalTok{(}\VariableTok{EbN0\_dB}\OperatorTok{/}\FloatTok{10}\NormalTok{)}\OperatorTok{;}

\CommentTok{\% Theoretical BER}
\VariableTok{BER\_theory} \OperatorTok{=} \VariableTok{qfunc}\NormalTok{(}\VariableTok{sqrt}\NormalTok{(}\FloatTok{2}\OperatorTok{*}\VariableTok{EbN0}\NormalTok{))}\OperatorTok{;}

\CommentTok{\% Simulated BER}
\VariableTok{N\_bits} \OperatorTok{=} \FloatTok{1e6}\OperatorTok{;}
\VariableTok{BER\_sim} \OperatorTok{=} \VariableTok{zeros}\NormalTok{(}\VariableTok{size}\NormalTok{(}\VariableTok{EbN0\_dB}\NormalTok{))}\OperatorTok{;}
\KeywordTok{for} \VariableTok{idx} \OperatorTok{=} \FloatTok{1}\OperatorTok{:}\VariableTok{length}\NormalTok{(}\VariableTok{EbN0}\NormalTok{)}
    \CommentTok{\% BPSK: transmit ±1}
    \VariableTok{bits} \OperatorTok{=} \FloatTok{2}\OperatorTok{*}\VariableTok{randi}\NormalTok{([}\FloatTok{0}\OperatorTok{,}\FloatTok{1}\NormalTok{]}\OperatorTok{,} \FloatTok{1}\OperatorTok{,} \VariableTok{N\_bits}\NormalTok{) }\OperatorTok{{-}} \FloatTok{1}\OperatorTok{;}
    \VariableTok{noise\_std} \OperatorTok{=} \FloatTok{1}\OperatorTok{/}\VariableTok{sqrt}\NormalTok{(}\FloatTok{2}\OperatorTok{*}\VariableTok{EbN0}\NormalTok{(}\VariableTok{idx}\NormalTok{))}\OperatorTok{;}
    \VariableTok{received} \OperatorTok{=} \VariableTok{bits} \OperatorTok{+} \VariableTok{noise\_std} \OperatorTok{*} \VariableTok{randn}\NormalTok{(}\FloatTok{1}\OperatorTok{,} \VariableTok{N\_bits}\NormalTok{)}\OperatorTok{;}
    \VariableTok{decisions} \OperatorTok{=} \VariableTok{sign}\NormalTok{(}\VariableTok{received}\NormalTok{)}\OperatorTok{;}
    \VariableTok{BER\_sim}\NormalTok{(}\VariableTok{idx}\NormalTok{) }\OperatorTok{=} \VariableTok{mean}\NormalTok{(}\VariableTok{decisions} \OperatorTok{\textasciitilde{}=} \VariableTok{bits}\NormalTok{)}\OperatorTok{;}
\KeywordTok{end}

\VariableTok{figure}\OperatorTok{;}
\VariableTok{semilogy}\NormalTok{(}\VariableTok{EbN0\_dB}\OperatorTok{,} \VariableTok{BER\_theory}\OperatorTok{,} \SpecialStringTok{\textquotesingle{}b{-}\textquotesingle{}}\OperatorTok{,} \SpecialStringTok{\textquotesingle{}LineWidth\textquotesingle{}}\OperatorTok{,} \FloatTok{2}\NormalTok{)}\OperatorTok{;}
\VariableTok{hold} \VariableTok{on}\OperatorTok{;}
\VariableTok{semilogy}\NormalTok{(}\VariableTok{EbN0\_dB}\OperatorTok{,} \VariableTok{BER\_sim}\OperatorTok{,} \SpecialStringTok{\textquotesingle{}ro\textquotesingle{}}\OperatorTok{,} \SpecialStringTok{\textquotesingle{}MarkerSize\textquotesingle{}}\OperatorTok{,} \FloatTok{6}\NormalTok{)}\OperatorTok{;}
\VariableTok{xlabel}\NormalTok{(}\SpecialStringTok{\textquotesingle{}E\_b/N\_0 (dB)\textquotesingle{}}\NormalTok{)}\OperatorTok{;} \VariableTok{ylabel}\NormalTok{(}\SpecialStringTok{\textquotesingle{}Bit Error Rate\textquotesingle{}}\NormalTok{)}\OperatorTok{;}
\VariableTok{title}\NormalTok{(}\SpecialStringTok{\textquotesingle{}BPSK BER: Theory vs Simulation\textquotesingle{}}\NormalTok{)}\OperatorTok{;}
\VariableTok{legend}\NormalTok{(}\SpecialStringTok{\textquotesingle{}Theory: Q(√(2E\_b/N\_0))\textquotesingle{}}\OperatorTok{,} \SpecialStringTok{\textquotesingle{}Simulation\textquotesingle{}}\NormalTok{)}\OperatorTok{;}
\VariableTok{grid} \VariableTok{on}\OperatorTok{;}
\end{Highlighting}
\end{Shaded}

\PSsection{3.10}{Gamma Distribution}

\PSsubsection{Physical Interpretation}

Models the \textbf{total waiting time until the $k-\mathrm{th}$ event} in a Poisson process, or any sum of $k$ independent exponential random variables.

\PSsubsection{Engineering Context}

\begin{itemize}
\tightlist
\item
  Total repair time for $k$ components (each with exponential repair time)
\item
  Time until $k-\mathrm{th}$ packet arrival
\item
  Aggregate service time for $k$ queued jobs
\item
  Rainfall amount modeling
\end{itemize}

\begin{PSbox}[unbreakable]{def}{Mathematical Definition}

$X \sim \mathrm{Gamma}(\alpha,\, \beta)$ where $\alpha$ = shape parameter, $\beta$ = scale parameter (or rate $\lambda = \frac{1}{\beta}$).

\end{PSbox}

\PSsubsection{PDF}

\PSdisplay{f_X(x) = \frac{x^{\alpha-1} e^{-x/\beta}}{\beta^\alpha \Gamma(\alpha)}, \quad x \geq 0}

where $\Gamma (\alpha) = (\alpha -1)$! for integer $\alpha$, and $\Gamma (\alpha) = \int_{0}^{\infty} t^{\alpha -1}e^{-t}dt$ in general.

\PSsubsection{CDF}

No closed form for general $\alpha$. Available via incomplete gamma function.

\PSsubsection{Mean}

\PSdisplay{E[X] = \alpha\beta}

\PSsubsection{Variance}

\PSdisplay{\text{Var}(X) = \alpha\beta^2}

\PSsubsection{Special Cases}

{\def\LTcaptype{none} % do not increment counter
\begin{longtable}[c]{@{}cc@{}}
\toprule\noalign{}
\rowcolor{PSnavy}\PSth{Parameters} & \PSth{Distribution} \\
\midrule\noalign{}
\endhead
\bottomrule\noalign{}
\endlastfoot
$\alpha = 1$ & $\mathrm{Exponential}(\beta)$ \\
$\alpha = \frac{n}{2},\, \beta = 2$ & Chi-square(n) \\
$\alpha$ = integer $k$ & $\mathrm{Erlang}(k,\, \beta)$ \\
\end{longtable}
}

\PSsubsection{When to Use}

\PSyes{} Sum of independent exponential lifetimes\PSbr{}\PSyes{} Waiting time for $k-\mathrm{th}$ Poisson event\PSbr{}\PSyes{} Skewed, non-negative data\PSbr{}\PSyes{} Flexible shape (adjustable skewness)

\PSsubsection{When NOT to Use}

\PSno{} Data can be negative\PSbr{}\PSno{} Data is symmetric (use Gaussian)\PSbr{}\PSno{} Simpler distribution fits (exponential for $k = 1$)

\PSsubsection{MATLAB Implementation}

\tcbset{PScodemode/.style={unbreakable}}

\begin{Shaded}
\begin{Highlighting}[]
\CommentTok{\%\% Gamma Distribution: Total Repair Time}
\CommentTok{\% 3 components must be repaired sequentially}
\CommentTok{\% Each repair time \textasciitilde{} Exp(mean = 2 hours)}
\VariableTok{alpha} \OperatorTok{=} \FloatTok{3}\OperatorTok{;}         \CommentTok{\% shape (number of stages)}
\VariableTok{beta} \OperatorTok{=} \FloatTok{2}\OperatorTok{;}          \CommentTok{\% scale (mean per stage)}

\VariableTok{x} \OperatorTok{=} \VariableTok{linspace}\NormalTok{(}\FloatTok{0}\OperatorTok{,} \FloatTok{20}\OperatorTok{,} \FloatTok{500}\NormalTok{)}\OperatorTok{;}
\VariableTok{pdf\_gamma} \OperatorTok{=} \VariableTok{gampdf}\NormalTok{(}\VariableTok{x}\OperatorTok{,} \VariableTok{alpha}\OperatorTok{,} \VariableTok{beta}\NormalTok{)}\OperatorTok{;}

\VariableTok{figure}\OperatorTok{;}
\VariableTok{plot}\NormalTok{(}\VariableTok{x}\OperatorTok{,} \VariableTok{pdf\_gamma}\OperatorTok{,} \SpecialStringTok{\textquotesingle{}b{-}\textquotesingle{}}\OperatorTok{,} \SpecialStringTok{\textquotesingle{}LineWidth\textquotesingle{}}\OperatorTok{,} \FloatTok{2}\NormalTok{)}\OperatorTok{;}
\VariableTok{xlabel}\NormalTok{(}\SpecialStringTok{\textquotesingle{}Total Repair Time (hours)\textquotesingle{}}\NormalTok{)}\OperatorTok{;}
\VariableTok{ylabel}\NormalTok{(}\SpecialStringTok{\textquotesingle{}f\_X(x)\textquotesingle{}}\NormalTok{)}\OperatorTok{;}
\VariableTok{title}\NormalTok{(}\VariableTok{sprintf}\NormalTok{(}\SpecialStringTok{\textquotesingle{}Gamma PDF: α=\%d, β=\%d\textquotesingle{}}\OperatorTok{,} \VariableTok{alpha}\OperatorTok{,} \VariableTok{beta}\NormalTok{))}\OperatorTok{;}
\VariableTok{grid} \VariableTok{on}\OperatorTok{;}

\CommentTok{\% Verify by summing exponentials}
\VariableTok{N} \OperatorTok{=} \FloatTok{100000}\OperatorTok{;}
\VariableTok{X\_sim} \OperatorTok{=} \VariableTok{sum}\NormalTok{(}\VariableTok{exprnd}\NormalTok{(}\VariableTok{beta}\OperatorTok{,} \VariableTok{alpha}\OperatorTok{,} \VariableTok{N}\NormalTok{)}\OperatorTok{,} \FloatTok{1}\NormalTok{)}\OperatorTok{;}  \CommentTok{\% Sum of 3 exponentials}
\VariableTok{fprintf}\NormalTok{(}\SpecialStringTok{\textquotesingle{}Mean: Theory=\%.2f, Sim=\%.2f\textbackslash{}n\textquotesingle{}}\OperatorTok{,} \VariableTok{alpha}\OperatorTok{*}\VariableTok{beta}\OperatorTok{,} \VariableTok{mean}\NormalTok{(}\VariableTok{X\_sim}\NormalTok{))}\OperatorTok{;}
\VariableTok{fprintf}\NormalTok{(}\SpecialStringTok{\textquotesingle{}Var:  Theory=\%.2f, Sim=\%.2f\textbackslash{}n\textquotesingle{}}\OperatorTok{,} \VariableTok{alpha}\OperatorTok{*}\VariableTok{beta}\OperatorTok{\^{}}\FloatTok{2}\OperatorTok{,} \VariableTok{var}\NormalTok{(}\VariableTok{X\_sim}\NormalTok{))}\OperatorTok{;}
\end{Highlighting}
\end{Shaded}

\PSsection{3.11}{Rayleigh Distribution}

\PSsubsection{Physical Interpretation}

Models the \textbf{magnitude (envelope)} of a $2D$ Gaussian random vector. When $X$ and $Y$ are independent $N(0,\, \sigma^{2})$, then $R = \sqrt{X^{2} + Y^{2}}$ follows a Rayleigh distribution.

\PSsubsection{Engineering Context}

\begin{itemize}
\tightlist
\item
  Envelope of narrowband noise (in-phase + quadrature)
\item
  Magnitude of wireless fading channel coefficient (Rayleigh fading)
\item
  Wind speed modeling
\item
  Distance error in $2D$ positioning (when errors are Gaussian in $x$ and $y$)
\end{itemize}

\begin{PSbox}[unbreakable]{def}{Mathematical Definition}

$R \sim \mathrm{Rayleigh}(\sigma)$, where $\sigma$ is the parameter (NOT the standard deviation of $R$).

\end{PSbox}

\PSsubsection{PDF}

\PSdisplay{f_R(r) = \frac{r}{\sigma^2} e^{-r^2/(2\sigma^2)}, \quad r \geq 0}

\PSsubsection{CDF}

\PSdisplay{F_R(r) = 1 - e^{-r^2/(2\sigma^2)}, \quad r \geq 0}

\PSsubsection{Mean}

\PSdisplay{E[R] = \sigma\sqrt{\frac{\pi}{2}} \approx 1.2533\sigma}

\PSsubsection{Variance}

\PSdisplay{\text{Var}(R) = \frac{4-\pi}{2}\sigma^2 \approx 0.4292\sigma^2}

\PSsubsection{Relationship to Other Distributions}

\begin{itemize}
\tightlist
\item
  $R^{2} \sim \mathrm{Exponential}(2\sigma^{2})$ — the power is exponentially distributed
\item
  If $X,\, Y \sim N(0,\, \sigma^{2})$ independent, then $\sqrt{X^{2}+Y^{2}} \sim \mathrm{Rayleigh}(\sigma)$
\end{itemize}

\PSsubsection{When to Use}

\PSyes{} Envelope/magnitude of complex Gaussian signal\PSbr{}\PSyes{} Rayleigh fading channel modeling (no line-of-sight)\PSbr{}\PSyes{} $2D$ distance with Gaussian errors in each dimension\PSbr{}\PSyes{} RSS (root-sum-square) of two independent Gaussian components

\PSsubsection{When NOT to Use}

\PSno{} Strong line-of-sight component present (use Rician)\PSbr{}\PSno{} Not modeling a magnitude/envelope\PSbr{}\PSno{} Components are not Gaussian or not equal-variance

\PSsubsection{MATLAB Implementation}

\tcbset{PScodemode/.style={unbreakable}}

\begin{Shaded}
\begin{Highlighting}[]
\CommentTok{\%\% Rayleigh Distribution: Wireless Fading Channel}
\VariableTok{sigma} \OperatorTok{=} \FloatTok{1}\OperatorTok{;}   \CommentTok{\% Rayleigh parameter}
\VariableTok{r} \OperatorTok{=} \VariableTok{linspace}\NormalTok{(}\FloatTok{0}\OperatorTok{,} \FloatTok{5}\OperatorTok{,} \FloatTok{500}\NormalTok{)}\OperatorTok{;}

\VariableTok{pdf\_ray} \OperatorTok{=} \VariableTok{raylpdf}\NormalTok{(}\VariableTok{r}\OperatorTok{,} \VariableTok{sigma}\NormalTok{)}\OperatorTok{;}
\VariableTok{cdf\_ray} \OperatorTok{=} \VariableTok{raylcdf}\NormalTok{(}\VariableTok{r}\OperatorTok{,} \VariableTok{sigma}\NormalTok{)}\OperatorTok{;}

\VariableTok{figure}\OperatorTok{;}
\VariableTok{subplot}\NormalTok{(}\FloatTok{2}\OperatorTok{,}\FloatTok{1}\OperatorTok{,}\FloatTok{1}\NormalTok{)}\OperatorTok{;}
\VariableTok{plot}\NormalTok{(}\VariableTok{r}\OperatorTok{,} \VariableTok{pdf\_ray}\OperatorTok{,} \SpecialStringTok{\textquotesingle{}b{-}\textquotesingle{}}\OperatorTok{,} \SpecialStringTok{\textquotesingle{}LineWidth\textquotesingle{}}\OperatorTok{,} \FloatTok{2}\NormalTok{)}\OperatorTok{;}
\VariableTok{xlabel}\NormalTok{(}\SpecialStringTok{\textquotesingle{}Channel Gain |h|\textquotesingle{}}\NormalTok{)}\OperatorTok{;} \VariableTok{ylabel}\NormalTok{(}\SpecialStringTok{\textquotesingle{}f\_R(r)\textquotesingle{}}\NormalTok{)}\OperatorTok{;}
\VariableTok{title}\NormalTok{(}\VariableTok{sprintf}\NormalTok{(}\SpecialStringTok{\textquotesingle{}Rayleigh PDF (σ = \%d)\textquotesingle{}}\OperatorTok{,} \VariableTok{sigma}\NormalTok{))}\OperatorTok{;}
\VariableTok{grid} \VariableTok{on}\OperatorTok{;}

\VariableTok{subplot}\NormalTok{(}\FloatTok{2}\OperatorTok{,}\FloatTok{1}\OperatorTok{,}\FloatTok{2}\NormalTok{)}\OperatorTok{;}
\VariableTok{plot}\NormalTok{(}\VariableTok{r}\OperatorTok{,} \VariableTok{cdf\_ray}\OperatorTok{,} \SpecialStringTok{\textquotesingle{}r{-}\textquotesingle{}}\OperatorTok{,} \SpecialStringTok{\textquotesingle{}LineWidth\textquotesingle{}}\OperatorTok{,} \FloatTok{2}\NormalTok{)}\OperatorTok{;}
\VariableTok{xlabel}\NormalTok{(}\SpecialStringTok{\textquotesingle{}r\textquotesingle{}}\NormalTok{)}\OperatorTok{;} \VariableTok{ylabel}\NormalTok{(}\SpecialStringTok{\textquotesingle{}P(R ≤ r)\textquotesingle{}}\NormalTok{)}\OperatorTok{;}
\VariableTok{title}\NormalTok{(}\SpecialStringTok{\textquotesingle{}Rayleigh CDF = Outage Probability\textquotesingle{}}\NormalTok{)}\OperatorTok{;}
\VariableTok{grid} \VariableTok{on}\OperatorTok{;}

\CommentTok{\% Outage probability: P(|h| \textless{} threshold)}
\VariableTok{threshold} \OperatorTok{=} \FloatTok{0.5}\OperatorTok{;}
\VariableTok{p\_outage} \OperatorTok{=} \VariableTok{raylcdf}\NormalTok{(}\VariableTok{threshold}\OperatorTok{,} \VariableTok{sigma}\NormalTok{)}\OperatorTok{;}
\VariableTok{fprintf}\NormalTok{(}\SpecialStringTok{\textquotesingle{}P(outage, |h| \textless{} \%.1f) = \%.4f\textbackslash{}n\textquotesingle{}}\OperatorTok{,} \VariableTok{threshold}\OperatorTok{,} \VariableTok{p\_outage}\NormalTok{)}\OperatorTok{;}
\end{Highlighting}
\end{Shaded}

\PSsubsection{Simulation Example: Rayleigh Fading}

\tcbset{PScodemode/.style={unbreakable}}

\begin{Shaded}
\begin{Highlighting}[]
\CommentTok{\%\% Generate Rayleigh Fading from Complex Gaussian}
\VariableTok{sigma} \OperatorTok{=} \FloatTok{1}\OperatorTok{;}
\VariableTok{N} \OperatorTok{=} \FloatTok{100000}\OperatorTok{;}

\CommentTok{\% Complex Gaussian channel coefficient}
\VariableTok{h} \OperatorTok{=} \VariableTok{sigma}\OperatorTok{/}\VariableTok{sqrt}\NormalTok{(}\FloatTok{2}\NormalTok{) }\OperatorTok{*}\NormalTok{ (}\VariableTok{randn}\NormalTok{(}\FloatTok{1}\OperatorTok{,}\VariableTok{N}\NormalTok{) }\OperatorTok{+} \FloatTok{1j}\OperatorTok{*}\VariableTok{randn}\NormalTok{(}\FloatTok{1}\OperatorTok{,}\VariableTok{N}\NormalTok{))}\OperatorTok{;}

\CommentTok{\% Magnitude is Rayleigh distributed}
\VariableTok{R} \OperatorTok{=} \VariableTok{abs}\NormalTok{(}\VariableTok{h}\NormalTok{)}\OperatorTok{;}

\VariableTok{figure}\OperatorTok{;}
\VariableTok{histogram}\NormalTok{(}\VariableTok{R}\OperatorTok{,} \FloatTok{100}\OperatorTok{,} \SpecialStringTok{\textquotesingle{}Normalization\textquotesingle{}}\OperatorTok{,} \SpecialStringTok{\textquotesingle{}pdf\textquotesingle{}}\NormalTok{)}\OperatorTok{;}
\VariableTok{hold} \VariableTok{on}\OperatorTok{;}
\VariableTok{r} \OperatorTok{=} \VariableTok{linspace}\NormalTok{(}\FloatTok{0}\OperatorTok{,} \FloatTok{5}\OperatorTok{,} \FloatTok{300}\NormalTok{)}\OperatorTok{;}
\VariableTok{plot}\NormalTok{(}\VariableTok{r}\OperatorTok{,} \VariableTok{raylpdf}\NormalTok{(}\VariableTok{r}\OperatorTok{,} \VariableTok{sigma}\OperatorTok{/}\VariableTok{sqrt}\NormalTok{(}\FloatTok{2}\NormalTok{))}\OperatorTok{,} \SpecialStringTok{\textquotesingle{}r{-}\textquotesingle{}}\OperatorTok{,} \SpecialStringTok{\textquotesingle{}LineWidth\textquotesingle{}}\OperatorTok{,} \FloatTok{2}\NormalTok{)}\OperatorTok{;}
\VariableTok{xlabel}\NormalTok{(}\SpecialStringTok{\textquotesingle{}|h|\textquotesingle{}}\NormalTok{)}\OperatorTok{;} \VariableTok{ylabel}\NormalTok{(}\SpecialStringTok{\textquotesingle{}PDF\textquotesingle{}}\NormalTok{)}\OperatorTok{;}
\VariableTok{title}\NormalTok{(}\SpecialStringTok{\textquotesingle{}Rayleigh Fading: Simulated vs Theory\textquotesingle{}}\NormalTok{)}\OperatorTok{;}
\VariableTok{legend}\NormalTok{(}\SpecialStringTok{\textquotesingle{}Simulation\textquotesingle{}}\OperatorTok{,} \SpecialStringTok{\textquotesingle{}Theory\textquotesingle{}}\NormalTok{)}\OperatorTok{;}
\end{Highlighting}
\end{Shaded}

\PSsection{3.12}{Chi-Square Distribution}

\PSsubsection{Physical Interpretation}

Models the \textbf{sum of squares of independent standard normal} random variables. Arises naturally in variance estimation and goodness-of-fit testing.

\PSsubsection{Engineering Context}

\begin{itemize}
\tightlist
\item
  Sample variance distribution (for Gaussian data)
\item
  Power of Gaussian noise (sum of squared components)
\item
  Goodness-of-fit test statistic
\item
  Confidence intervals for variance
\item
  Energy detection in spectrum sensing
\end{itemize}

\begin{PSbox}[unbreakable]{def}{Mathematical Definition}

If $Z_{1},\, Z_{2},\, \ldots,\, Z_{n}$ are independent $N(0,\,1)$, then:\PSdisplay{\chi^2 = Z_1^2 + Z_2^2 + \cdots + Z_n^2 \sim \chi^2(n)}

$n$ = degrees of freedom.

\end{PSbox}

\PSsubsection{PDF}

\PSdisplay{f_X(x) = \frac{x^{n/2-1} e^{-x/2}}{2^{n/2}\Gamma(n/2)}, \quad x \geq 0}

\PSsubsection{CDF}

No closed form for general $n$. Computed via incomplete gamma function.

\PSsubsection{Mean}

\PSdisplay{E[X] = n}

\PSsubsection{Variance}

\PSdisplay{\text{Var}(X) = 2n}

\PSsubsection{Relationship to Other Distributions}

\begin{itemize}
\tightlist
\item
  Chi-square(n) $= \mathrm{Gamma}\left(\frac{n}{2},\, 2\right)$
\item
  Chi-square(1) = square of $N(0,\,1)$
\item
  Chi-square(2) $= \mathrm{Exponential}\left(\frac{1}{2}\right)$
\item
  For large $n$: $\chi^{2}(n) \approx N(n,\, 2n)$ by CLT
\end{itemize}

\PSsubsection{When to Use}

\PSyes{} Sum of squared Gaussian random variables\PSbr{}\PSyes{} Variance testing and confidence intervals\PSbr{}\PSyes{} Goodness-of-fit tests\PSbr{}\PSyes{} Energy detection in $N$ dimensions

\PSsubsection{When NOT to Use}

\PSno{} Variables are not Gaussian\PSbr{}\PSno{} Variables are not independent\PSbr{}\PSno{} Variables are not standard normal (need scaling)

\PSsubsection{MATLAB Implementation}

\tcbset{PScodemode/.style={unbreakable}}

\begin{Shaded}
\begin{Highlighting}[]
\CommentTok{\%\% Chi{-}Square Distribution: Noise Power}
\CommentTok{\% Noise power from N complex dimensions = sum of 2N squared Gaussians}

\VariableTok{n\_dof} \OperatorTok{=} \FloatTok{10}\OperatorTok{;}        \CommentTok{\% degrees of freedom}
\VariableTok{x} \OperatorTok{=} \VariableTok{linspace}\NormalTok{(}\FloatTok{0}\OperatorTok{,} \FloatTok{30}\OperatorTok{,} \FloatTok{500}\NormalTok{)}\OperatorTok{;}

\VariableTok{pdf\_chi2} \OperatorTok{=} \VariableTok{chi2pdf}\NormalTok{(}\VariableTok{x}\OperatorTok{,} \VariableTok{n\_dof}\NormalTok{)}\OperatorTok{;}
\VariableTok{cdf\_chi2} \OperatorTok{=} \VariableTok{chi2cdf}\NormalTok{(}\VariableTok{x}\OperatorTok{,} \VariableTok{n\_dof}\NormalTok{)}\OperatorTok{;}

\VariableTok{figure}\OperatorTok{;}
\VariableTok{plot}\NormalTok{(}\VariableTok{x}\OperatorTok{,} \VariableTok{pdf\_chi2}\OperatorTok{,} \SpecialStringTok{\textquotesingle{}b{-}\textquotesingle{}}\OperatorTok{,} \SpecialStringTok{\textquotesingle{}LineWidth\textquotesingle{}}\OperatorTok{,} \FloatTok{2}\NormalTok{)}\OperatorTok{;}
\VariableTok{xlabel}\NormalTok{(}\SpecialStringTok{\textquotesingle{}χ² Value (Noise Power)\textquotesingle{}}\NormalTok{)}\OperatorTok{;}
\VariableTok{ylabel}\NormalTok{(}\SpecialStringTok{\textquotesingle{}f\_X(x)\textquotesingle{}}\NormalTok{)}\OperatorTok{;}
\VariableTok{title}\NormalTok{(}\VariableTok{sprintf}\NormalTok{(}\SpecialStringTok{\textquotesingle{}Chi{-}Square PDF: \%d degrees of freedom\textquotesingle{}}\OperatorTok{,} \VariableTok{n\_dof}\NormalTok{))}\OperatorTok{;}
\VariableTok{grid} \VariableTok{on}\OperatorTok{;}

\CommentTok{\% Simulation: sum of squared standard normals}
\VariableTok{N} \OperatorTok{=} \FloatTok{100000}\OperatorTok{;}
\VariableTok{Z} \OperatorTok{=} \VariableTok{randn}\NormalTok{(}\VariableTok{n\_dof}\OperatorTok{,} \VariableTok{N}\NormalTok{)}\OperatorTok{;}
\VariableTok{X\_sim} \OperatorTok{=} \VariableTok{sum}\NormalTok{(}\VariableTok{Z}\OperatorTok{.\^{}}\FloatTok{2}\OperatorTok{,} \FloatTok{1}\NormalTok{)}\OperatorTok{;}

\VariableTok{fprintf}\NormalTok{(}\SpecialStringTok{\textquotesingle{}Mean: Theory=\%d, Sim=\%.2f\textbackslash{}n\textquotesingle{}}\OperatorTok{,} \VariableTok{n\_dof}\OperatorTok{,} \VariableTok{mean}\NormalTok{(}\VariableTok{X\_sim}\NormalTok{))}\OperatorTok{;}
\VariableTok{fprintf}\NormalTok{(}\SpecialStringTok{\textquotesingle{}Var:  Theory=\%d, Sim=\%.2f\textbackslash{}n\textquotesingle{}}\OperatorTok{,} \FloatTok{2}\OperatorTok{*}\VariableTok{n\_dof}\OperatorTok{,} \VariableTok{var}\NormalTok{(}\VariableTok{X\_sim}\NormalTok{))}\OperatorTok{;}
\end{Highlighting}
\end{Shaded}

\PSsection{3.13}{Distribution Selection Guide}

\PSsubsection{Decision Tree: Discrete Distributions}

\PSfigure{m03-tree-discrete}

\PSsubsection{Decision Tree: Continuous Distributions}

\PSfigure{m03-tree-continuous}

\PSsubsection{Quick Recognition Patterns}

{\def\LTcaptype{none} % do not increment counter
\begin{longtable}[c]{@{}cc@{}}
\toprule\noalign{}
\rowcolor{PSnavy}\PSth{If you see…} & \PSth{Think…} \\
\midrule\noalign{}
\endhead
\bottomrule\noalign{}
\endlastfoot
Binary outcome & Bernoulli \\
``$n$ trials, $k$ successes'' & Binomial \\
``How many until first…'' & Geometric \\
``Events per interval'' & Poisson \\
``Equal chance anywhere in $[a,\,b]$'' & Uniform \\
``Time until next event'' & Exponential \\
``Sum of many effects'' or ``noise'' & Gaussian \\
``Total waiting time for $k$ events'' & Gamma \\
``Signal envelope'' or ``fading'' & Rayleigh \\
``Sum of squared normals'' & Chi-square \\
\end{longtable}
}

\PSsection{3.14}{Summary Table}

{\def\LTcaptype{none} % do not increment counter
\begin{longtable}[c]{@{}
  >{\centering\arraybackslash}p{(0.965\linewidth - 10\tabcolsep) * \real{0.2407}}
  >{\centering\arraybackslash}p{(0.965\linewidth - 10\tabcolsep) * \real{0.1111}}
  >{\centering\arraybackslash}p{(0.965\linewidth - 10\tabcolsep) * \real{0.2037}}
  >{\centering\arraybackslash}p{(0.965\linewidth - 10\tabcolsep) * \real{0.1111}}
  >{\centering\arraybackslash}p{(0.965\linewidth - 10\tabcolsep) * \real{0.1852}}
  >{\centering\arraybackslash}p{(0.965\linewidth - 10\tabcolsep) * \real{0.1481}}@{}}
\toprule\noalign{}
\rowcolor{PSnavy}\PSth{Distribution} & \PSth{Type} & \PSth{Parameters} & \PSth{Mean} & \PSth{Variance} & \PSth{MATLAB} \\
\midrule\noalign{}
\endhead
\bottomrule\noalign{}
\endlastfoot
Bernoulli & Discrete & $p$ & $p$ & $p(1-p)$ & \PScode{binornd(1,p)} \\
Binomial & Discrete & $n,\, p$ & $np$ & $np(1-p)$ & \PScode{binornd(n,p)} \\
Geometric & Discrete & $p$ & $\frac{1}{p}$ & $\frac{1-p}{p^{2}}$ & \PScode{geornd(p)+1} \\
Neg. Binomial & Discrete & $r,\, p$ & $\frac{r}{p}$ & $\frac{r(1-p)}{p^{2}}$ & \PScode{nbinrnd(r,p)+r} \\
Poisson & Discrete & $\lambda$ & $\lambda$ & $\lambda$ & \PScode{poissrnd(λ)} \\
Uniform & Continuous & a, $b$ & $\frac{a+b}{2}$ & $\frac{(b-a)^{2}}{12}$ & \PScode{unifrnd(a,b)} \\
Exponential & Continuous & $\lambda$ (or $\beta = \frac{1}{\lambda}$) & $\frac{1}{\lambda}$ & $\frac{1}{\lambda^{2}}$ & \PScode{exprnd(β)} \\
Gaussian & Continuous & $\mu,\, \sigma^{2}$ & $\mu$ & $\sigma^{2}$ & \PScode{normrnd(μ,σ)} \\
Gamma & Continuous & $\alpha,\, \beta$ & $\alpha \beta$ & $\alpha \beta^{2}$ & \PScode{gamrnd(α,β)} \\
Rayleigh & Continuous & $\sigma$ & $\sigma \sqrt{\frac{\pi}{2}}$ & $\frac{(4-\pi)\sigma^{2}}{2}$ & \PScode{raylrnd(σ)} \\
Chi-square & Continuous & $n$ & $n$ & $2n$ & \PScode{chi2rnd(n)} \\
\end{longtable}
}

\PSsection{3.15}{Practice Problems}

\begin{PSbox}[unbreakable]{prob}{Problem 1: Distribution Identification}

Identify the appropriate distribution for each scenario:

\begin{enumerate}
\def\labelenumi{(\alph{enumi})}
\tightlist
\item
  Number of dropped calls in a 10-minute window (average: 2 per 10 min)
\item
  Voltage of thermal noise across a resistor
\item
  Time until a server crashes (constant failure rate)
\item
  Number of defective ICs out of a batch of 100 (defect rate 3\%)
\item
  Envelope of received signal through multipath with no LOS
\end{enumerate}

\PSsolution{} (a) $\mathrm{Poisson}(\lambda = 2)$, (b) $\mathrm{Gaussian}(0,\, \sigma^{2})$, (c) $\mathrm{Exponential}(\lambda)$, (d) $\mathrm{Binomial}(100,\, 0.03)$, (e) $\mathrm{Rayleigh}(\sigma)$

\end{PSbox}

\begin{PSbox}[unbreakable]{prob}{Problem 2: Exponential vs Poisson}

Packets arrive at a router as a Poisson process with rate $\lambda = 5$ packets/second.

\begin{enumerate}
\def\labelenumi{(\alph{enumi})}
\tightlist
\item
  What is $P(\text{more than 8 packets in 1 second})$?
\item
  What is $P(\text{no packet arrives in the next 0.5 seconds})$?
\item
  What distribution describes the inter-arrival time?
\item
  What is the mean inter-arrival time?
\end{enumerate}

\PSsolution{} 

\begin{enumerate}
\def\labelenumi{(\alph{enumi})}
\tightlist
\item
  $X \sim \mathrm{Poisson}(5)$: $P(X > 8) = 1 - \texttt{poisscdf}(8,\, 5) = 0.0681$
\item
  $T \sim \mathrm{Exp}(5)$: $P(T > 0.5) = e^{-5 \times 0.5} = e^{-2.5} = 0.0821$
\item
  Exponential with rate $\lambda = 5$
\item
  Mean $= \frac{1}{\lambda} = 0.2$ seconds
\end{enumerate}

\end{PSbox}

\begin{PSbox}[unbreakable]{prob}{Problem 3: Gaussian Probability}

A resistor has nominal value 1000\ensuremath{\Omega} with manufacturing variation modeled as $N(1000,\, 25) (\sigma = 5\,\Omega)$.

\begin{enumerate}
\def\labelenumi{(\alph{enumi})}
\tightlist
\item
  What fraction of resistors are within $\pm 10\,\Omega$ of nominal (spec: $990-1010\,\Omega$)?
\item
  What fraction fail the $\pm 1\%$ tolerance test?
\item
  What tolerance range contains 99.7\% of resistors?
\end{enumerate}

\PSsolution{} 

\begin{enumerate}
\def\labelenumi{(\alph{enumi})}
\tightlist
\item
  $P(990 \le X \le 1010) = P(-2 \le Z \le 2) = 0.9545 \to 95.45\%$
\item
  $\pm 1\% = \pm 10\,\Omega$ \ensuremath{\to} same as (a), fail rate $= 1 - 0.9545 = 4.55\%$
\item
  $\pm 3\sigma = \pm 15\,\Omega$ \ensuremath{\to} range $[985,\, 1015]\Omega$
\end{enumerate}

\emph{Next Module: {Module 4 — Expectation and Moments}}

\end{PSbox}
